% Nouveau EGA V (cote 25): archive cote 25, site page 39 (PDF page 40).
% typescript. Draft compared: none.
% AI-transcribed by Claude Opus 5.5 (claude-opus-5-5, adaptive thinking, effort high), Message Batches API, batch ega25b (msgbatch_01A7n5YuXubUPNK1zsepi8RG), 2026-11-13; not yet checked by eye.
\scanpage{25/39}
\mspage{\hand{\unsure{35|33}}}

une sous-variété linéaire passant par a son espace tangent en a, consi-%
déré comme sous-espace \var{\xout{tangent}}{}{typewriter strike-out}{ega25_25_039-v0001} de l'espace tangent à P en a. Ce genre
\var{\margin{\hand{||}}}{\placedmargin{\hand{||}}}{marginal note placed}{ega25_25_039-v0002}de \unsure{sorites|sorttes}, ensemble avec divers sorites sur les sous-variétés linéaires
et les grassmaniennes, devraient être donnés dans un ou deux n$^{\text{o}}$s prélimi
naires, bien entendu en énonçant tout systématiquement sur une base
quelconque.

En fait, il y a mieux, savoir que le \uline{sous-préschéma} $\underline{Y}^{\mathrm{sing}}$
des zéros singuliers de $\hand{\phi}$ relativement à $\overset{\hand{\vee}}{P}$ , défini dans par.16, n'est
autre que $P(\hand{\unsure{\underline{\nu}|\check{\underline{\nu}}}}_{X/P})$ , et (comme ce dernier est lisse sur S de dimension
relative d \unsure{$+$} (r-d-1) = r - 1 (r étant la dim relative de P sur S)
on est sous les conditions fav\unsure{o|n}rables étudiées au par.16 . \var{\deleted{Pour vérifi
er cette assertion, on utilise le fait suivant (qui devrait figurer au
par. 17 des morphismes lisses), en l'appliquant à $X_{\overset{\hand{\vee}}{P}}$ \var{\inserted[typed]{et $\underline{P}^1(\underline{O}\ (1,1))$}}{et $\underline{P}^1(\underline{O}\ (1,1))$}{insertion}{ega25_25_039-v0004} S : Soit
X lisse sur S de dimension relative m, $\underline{M}$ un Module localement libre sur
X de rang n, $\hand{\mu}$ une section de $\underline{M}$ , x V( $\hand{\mu}$ )=Z , supposons que
dim$_x$ Z$_s$ = n-m , \var{\xout{et}}{}{typewriter strike-out}{ega25_25_039-v0005} (resp. et que $Z_{\overline{s}\ \text{réd}}$ soit lisse sur $\overline{k(s)}$ en $\bar{x}$ )
alors Z est \var{\xout{lisse sur}}{}{typewriter strike-out}{ega25_25_039-v0006} plat sur S en x (resp. lisse sur S en x). (Dém:
on peut supposer M = $\underline{O}_X^{\,m}$ , donc $\hand{\mu} = (\hand{\mu}_1,..,\hand{\mu}_n)$ , alors il résulte de
la théorie de la dimension et du fait que $\underline{O}_{X_s,x}$ est CM que les
$\hand{\mu}_{i\ s}$ forment une suite $\underline{O}_{X_s}$-régulière en x, donc \var{\deleted{les $\hand{\mu}_i$ sont transvers
salement $\underline{O}_X$-réguliers rel. à S; di de plus Z}}{}{deletion}{ega25_25_039-v0007}}}{}{deletion}{ega25_25_039-v0003} Pour le vérifier,
notons que par définition, $\underline{Y}^{\mathrm{sing}}$ n'est autre que le sous-préschéma \var{\inserted[typed]{de $\underline{Y}$}}{de $\underline{Y}$}{insertion}{ega25_25_039-v0008}
des zéros de la section \var{\inserted{$\psi =$}}{$\psi =$}{insertion}{ega25_25_039-v0009} $d\hand{\phi}|\underline{Y}$ de $\hand{\underline{\Omega}}^1_{X_{\overset{\hand{\vee}}{P}}/\overset{\hand{\vee}}{P}}(1,1)\otimes \underline{O}_Y = \hand{\underline{\Omega}}^1_{X/S}\otimes\underline{O}_Y(1,1)$\hand{.}
Nous allons donner une autre interprétation de cette section, dont
la conclusion résultera aussitôt. Pour ceci, considérons le diagramme
\var{\deleted{de Modules}}{}{deletion}{ega25_25_039-v0010} \var{\inserted{suites exactes}}{suites exactes}{insertion}{ega25_25_039-v0011}
suivant sur $X_{\overset{\hand{\vee}}{P}}$ , ou plus généralement sur un préschéma Z \var{\inserted{quelconque}}{quelconque}{insertion}{ega25_25_039-v0012} au dessus de
$X_{\overset{\hand{\vee}}{P}}$ :

\mspage{\hand{39}}
